\documentclass[12pt]{article}
\usepackage{mathptmx}
\usepackage[utf8]{inputenc}
\usepackage{cmap}
\usepackage{mathtext}
\usepackage{array}
\usepackage{booktabs}
\usepackage{multirow}
\usepackage[russian,english]{babel}
\usepackage{amsmath,amsfonts,amssymb,amsthm,mathtools} % AMS
\usepackage{icomma}
\usepackage[T2A]{fontenc}
\usepackage[left=10mm, top=20mm, right=18mm, bottom=15mm, footskip=10mm]{geometry}
\usepackage{indentfirst}
\usepackage{euscript}	 % Шрифт Евклид
\usepackage{mathrsfs} % Красивый матшрифт
\usepackage{amsmath,amssymb}
\usepackage[italicdiff]{physics}
\usepackage{graphicx}
\graphicspath{{images/}}
\DeclareGraphicsExtensions{.pdf,.png,.jpg}
\usepackage{wrapfig}
\usepackage{subcaption}
\usepackage{caption}
\usepackage{hyperref}
\usepackage[rgb]{xcolor}
\hypersetup{				% Гиперссылки
    colorlinks=true,       	% false: ссылки в рамках
	urlcolor=blue          % на URL
}
\captionsetup[figure]{name=Рисунок}
\captionsetup[table]{name=Таблица}


% Определяем команду для строки оглавления с отточками
\newcommand{\tocline}[2]{%
  \noindent\makebox[0pt][l]{#1}\hfill
  \leaders\hbox to 1.5mm{\hss.\hss}\hfill
  \makebox[0pt][r]{#2}\par
}

% Альтернативный стиль (более надёжный, работает всегда):
\newcommand{\tocentry}[2]{%
  \noindent
  \parbox[t]{0.8\textwidth}{\raggedright #1}%
  \parbox[t]{0.15\textwidth}{\raggedleft #2}%
  \par
  \vspace{-1.2\baselineskip}
  \noindent\leaders\hbox to 1.5mm{\hss.\hss}\hfill\mbox{}%
  \par
  \vspace{0.8\baselineskip}
}

% Но самый простой и надёжный способ — через \dotfill:
\newcommand{\tocitem}[2]{%
  \noindent #1 \dotfill\ #2\par
}


\begin{document}

\thispagestyle{empty}

\begin{center}

% Верхняя часть титульного листа
\large
Министерство науки и высшего образования Российской Федерации \\
\large
ФЕДЕРАЛЬНОЕ ГОСУДАРСТВЕННОЕ АВТОНОМНОЕ \\
ОБРАЗОВАТЕЛЬНОЕ УЧРЕЖДЕНИЕ ВЫСШЕГО ОБРАЗОВАНИЯ \\
\large
«МОСКОВСКИЙ ФИЗИКО-ТЕХНИЧЕСКИЙ ИНСТИТУТ \\
(НАЦИОНАЛЬНЫЙ ИССЛЕДОВАТЕЛЬСКИЙ УНИВЕРСИТЕТ)» \\
(МФТИ, Физтех)

\vspace{4cm}

% Основная информация
\large
КАФЕДРА ЭЛЕКТРОНИКИ

\vspace{0.7cm}

\Large
ОТЧЕТ \\
ПО ЛАБОРАТОРНОЙ РАБОТЕ

\vspace{0.7cm}

\large
\textbf{Конвективная диффузия 
в молекулярно-электронных  
преобразователях }

\vfill

% Блок с подписями
\noindent
Работу выполнил \hfill \underline{\hspace{6cm}} Д.А.Воронин Б04-407 \\
\hfill \small{(подпись, дата)}\par

\vspace{3cm}

\noindent
Работу принял, оценка \hfill \underline{\hspace{6cm}} \\
\hfill \small{(подпись, дата, оценка)}

\vfill

% Нижняя часть
Долгопрудный 2026

\end{center}


%%%%%%%%%%%%%

\thispagestyle{empty}

\begin{center}

% Верхняя часть титульного листа
\large
Министерство науки и высшего образования Российской Федерации \\
\large
ФЕДЕРАЛЬНОЕ ГОСУДАРСТВЕННОЕ АВТОНОМНОЕ \\
ОБРАЗОВАТЕЛЬНОЕ УЧРЕЖДЕНИЕ ВЫСШЕГО ОБРАЗОВАНИЯ \\
\large
«МОСКОВСКИЙ ФИЗИКО-ТЕХНИЧЕСКИЙ ИНСТИТУТ \\
(НАЦИОНАЛЬНЫЙ ИССЛЕДОВАТЕЛЬСКИЙ УНИВЕРСИТЕТ)» \\
(МФТИ, Физтех)

\vspace{4cm}

% Основная информация
\large
КАФЕДРА ЭЛЕКТРОНИКИ

\vspace{0.7cm}

\Large
ОТЧЕТ \\
ПО ЛАБОРАТОРНОЙ РАБОТЕ

\vspace{0.7cm}

\large
\textbf{Конвективная диффузия 
в молекулярно-электронных  
преобразователях }

\vfill

% Блок с подписями
\noindent
Работу выполнил \hfill \underline{\hspace{6cm}} Н.В.Дрожжин Б04-407 \\
\hfill \small{(подпись, дата)}\par

\vspace{3cm}

\noindent
Работу принял, оценка \hfill \underline{\hspace{6cm}} \\
\hfill \small{(подпись, дата, оценка)}

\vfill

% Нижняя часть
Долгопрудный 2026

\end{center}


%%%%%%%%%%%%%%

\thispagestyle{empty}

\begin{center}

% Верхняя часть титульного листа
\large
Министерство науки и высшего образования Российской Федерации \\
\large
ФЕДЕРАЛЬНОЕ ГОСУДАРСТВЕННОЕ АВТОНОМНОЕ \\
ОБРАЗОВАТЕЛЬНОЕ УЧРЕЖДЕНИЕ ВЫСШЕГО ОБРАЗОВАНИЯ \\
\large
«МОСКОВСКИЙ ФИЗИКО-ТЕХНИЧЕСКИЙ ИНСТИТУТ \\
(НАЦИОНАЛЬНЫЙ ИССЛЕДОВАТЕЛЬСКИЙ УНИВЕРСИТЕТ)» \\
(МФТИ, Физтех)

\vspace{4cm}

% Основная информация
\large
КАФЕДРА ЭЛЕКТРОНИКИ

\vspace{0.7cm}

\Large
ОТЧЕТ \\
ПО ЛАБОРАТОРНОЙ РАБОТЕ

\vspace{0.7cm}

\large
\textbf{ Конвективная диффузия 
в молекулярно-электронных  
преобразователях }

\vfill

% Блок с подписями
\noindent
Работу выполнил \hfill \underline{\hspace{6cm}} Т.А.Ганицев Б04-407 \\
\hfill \small{(подпись, дата)}\par

\vspace{3cm}

\noindent
Работу принял, оценка \hfill \underline{\hspace{6cm}} \\
\hfill \small{(подпись, дата, оценка)}

\vfill

% Нижняя часть
Долгопрудный 2026

\end{center}

%%%%%%%%%%%

\thispagestyle{empty}

\begin{center}

% Верхняя часть титульного листа
\large
Министерство науки и высшего образования Российской Федерации \\
\large
ФЕДЕРАЛЬНОЕ ГОСУДАРСТВЕННОЕ АВТОНОМНОЕ \\
ОБРАЗОВАТЕЛЬНОЕ УЧРЕЖДЕНИЕ ВЫСШЕГО ОБРАЗОВАНИЯ \\
\large
«МОСКОВСКИЙ ФИЗИКО-ТЕХНИЧЕСКИЙ ИНСТИТУТ \\
(НАЦИОНАЛЬНЫЙ ИССЛЕДОВАТЕЛЬСКИЙ УНИВЕРСИТЕТ)» \\
(МФТИ, Физтех)

\vspace{4cm}

% Основная информация
\large
КАФЕДРА ЭЛЕКТРОНИКИ

\vspace{0.7cm}

\Large
ОТЧЕТ \\
ПО ЛАБОРАТОРНОЙ РАБОТЕ

\vspace{0.7cm}

\large
\textbf{Конвективная диффузия 
в молекулярно-электронных  
преобразователях }

\vfill

% Блок с подписями
\noindent
Работу выполнил \hfill \underline{\hspace{6cm}} В.С.Анисимов Б04-407 \\
\hfill \small{(подпись, дата)}\par

\vspace{3cm}

\noindent
Работу принял, оценка \hfill \underline{\hspace{6cm}} \\
\hfill \small{(подпись, дата, оценка)}

\vfill

% Нижняя часть
Долгопрудный 2026

\end{center}

\end{document}
